\documentclass[10pt,a4paper,sans]{moderncv}

\moderncvstyle{banking}
\moderncvcolor{blue}
\usepackage[scale=0.90]{geometry}
\nopagenumbers{}

\name{Andrés}{Pescara Montero}
\title{Data Engineer}
\address{Viña del Mar, Chile}{}{}
\phone[mobile]{+56 990026042}
\email{andrespescara@gmail.com}
\social[linkedin]{apescaram}
\social[github]{apescara}

\begin{document}

\makecvtitle

\section{Perfil profesional}
Data Engineer con más de siete años desarrollando y productivizando soluciones de datos y Machine Learning. Diseña arquitecturas escalables en GCP, integra fuentes batch, APIs y streaming, y automatiza pipelines con Python, SQL, Airflow y dbt. Experiencia en CI/CD, gobierno, eficiencia de costos, liderazgo técnico y colaboración con Data Science y negocio.

\section{Experiencia}
\cventry{jul. 2025--actualidad}{Machine Learning Engineer}{Sodimac}{Chile}{}{
  \begin{itemize}
    \item Redujo de siete a dos minutos la duración promedio del pipeline de CI en Cloud Build mediante imágenes Docker base, Bash, uv y una actualización controlada de Python.
    \item Administra con Terraform Cloud Composer, GKE y Vertex AI Workbench en una plataforma de cerca de 200 repositorios, 260 DAGs de Airflow y más de 40 APIs.
    \item Estandarizó el feature engineering con dbt, implementó selección dinámica de máquinas según consumo histórico y gestiona el backlog y paneles de gasto del equipo de Data Engineering.
  \end{itemize}
}

\cventry{nov. 2023--jun. 2025}{Data Engineer}{The Wild Brands}{Chile y México}{}{
  \begin{itemize}
    \item Diseñó desde cero un Data Lake en GCP para cerca de 150 usuarios, con ambientes productivos y de prueba, capas raw--stage--share y un modelo dbt de más de 200 entidades.
    \item Integró SAP HANA, Odoo, Airbyte y Shopify; combinó procesos batch con webhooks procesados mediante Pub/Sub y Dataflow bajo un modelo de datos común.
    \item Gestionó la infraestructura con Terraform, desarrolló pipelines de GitHub Actions y lideró integrantes junior a medida que creció el equipo.
  \end{itemize}
}

\cventry{jul. 2023--oct. 2023}{Senior Data Engineer}{Falabella Tecnologías Corporativas}{Chile}{}{
  \begin{itemize}
    \item Mantuvo pipelines de datos y ML en Python, Docker, Cloud Composer, GKE, BigQuery y Cloud Storage para múltiples equipos de Data Science; asumió interinamente el liderazgo técnico.
  \end{itemize}
}

\cventry{mar. 2021--jul. 2023}{Senior Data Engineer}{Sodimac}{Chile}{}{
  \begin{itemize}
    \item Diseñó una plataforma de ML con Docker, Cloud Composer y GKE; migró Airflow 1 a 2 sin interrupciones y creó operadores para prevenir fallas por cuotas.
    \item Implementó CI/CD con GitLab CI y creó en BigQuery un repositorio de features que integraba APIs, web scraping e indicadores externos.
  \end{itemize}
}

\cventry{jun. 2019--feb. 2021}{Consultor}{Management Solutions}{Chile}{}{
  \begin{itemize}
    \item Desarrolló ETLs y controles de calidad para banca con SQL, DataStage, Control-M y DB2; implementó pipelines regulatorios bronze--silver--gold en SQL Server y una prueba de concepto en BigQuery.
  \end{itemize}
}

\section{Competencias técnicas}
\cvitem{Datos}{Python, SQL, Airflow/Cloud Composer, BigQuery, dbt, Dataflow, Pub/Sub, Airbyte, modelado y calidad de datos}
\cvitem{Cloud y DevOps}{GCP, Docker, Kubernetes/GKE, Terraform, Cloud Build, GitLab CI, GitHub Actions, Git, Cloud Functions, Cloud Run}
\cvitem{Prácticas}{Arquitectura de Data Lake, CI/CD, FinOps, MLOps, liderazgo técnico y mentoría}

\section{Educación y certificaciones}
\cventry{2014--2019}{Ingeniería Civil Informática}{Universidad Adolfo Ibáñez}{}{}{}
\cvitem{2025--2027}{Professional Machine Learning Engineer — Google Cloud}
\cvitem{2019}{Professional Data Engineer — Google Cloud (vigencia finalizada en 2021)}

\section{Idiomas}
\cvitem{}{Español nativo; inglés avanzado (First Certificate in English, Cambridge, 2012)}

\end{document}
